\section{Discussion and Failure Cases}

Building the system surfaced several instructive failure modes, most of them in the local
layer.

\paragraph{Dead ends (local planning alone).} As argued in Section~1, a local optimiser
walks into U-shaped obstacles because escaping requires temporarily increasing the
goal distance. My global A* layer removes this failure by planning a complete route the
local planner only has to track.

\paragraph{DWA rotate-in-place / frozen robot.} DWA fails structurally in clutter (0/25).
From a standstill the acceleration-limited window admits only near-zero speeds, so a
forward candidate barely displaces the robot over the horizon; the progress terms cannot
distinguish candidates while the obstacle term can, and pivoting in place wins. The robot
never builds speed, so the window never opens resulting into a frozen-robot local
minimum. A related earlier bug amplified this: querying the risk field by nearest-cell
rounding created cost discontinuities at cell boundaries; replacing it with bilinear
interpolation removed those artefacts. DWA also rejects any rollout that clips an inflated
cell as infinite cost, which near corners can eliminate every forward option.

\paragraph{Why MPPI helps, and where it still fails.} MPPI avoids these traps by optimising
full control sequences and softly averaging over rollouts, reaching 68\% of goals. Its
remaining 32\% of failures occur in the tightest passages, where the sampling distribution
rarely proposes the narrow feasible manoeuvre within the horizon; a longer horizon or more
samples trades compute for a higher success rate. In these 32% there are some paths which are simply cannot be executed because of the map layout.

\paragraph{Trade-offs.} Across both planners the recurring tension is speed vs.\ safety
vs.\ efficiency: cautious weightings are safe but slow and freeze-prone, aggressive
weightings are fast but graze obstacles. The probabilistic field mitigates this by turning
a hard constraint into a smooth gradient, but the weights still require fine-tuning.
